\section{Nonrigid classes on a quotient resolution}
\label{sec:kummer-boundary-classes}

An irrational nef boundary ray can admit many smooth semipositive representatives. A transverse periodic function produces them on a torus. Making the functions vanish near the fixed points of the involution allows the forms to pass smoothly to its quotient resolution.

Let $A=\mathbb C^2/\Lambda$ be a complex two-torus, let $Y=A/\{\pm1\}$, and let $r:X\to Y$ be the minimal resolution. The sixteen singular points of $Y$ are the images of $A[2]=(\frac12\Lambda)/\Lambda$.

\Needspace{12\baselineskip}
\begin{lemma}\label{lem:kummer-resolution-kahler}\index{Kummer resolution}
The surface $X$ is compact K\"ahler. Its exceptional set consists of sixteen disjoint projective lines with self-intersection $-2$.
\end{lemma}
\begin{proof}
Blow up the sixteen fixed points on $A$ and take the quotient by the lifted involution. Near one exceptional line, a blowup chart has $z_1=u$, $z_2=uv$, and the involution is $(u,v)\mapsto(-u,v)$. The quotient chart has smooth coordinates $(s,v)=(u^2,v)$. On the other chart, $v'=1/v$ and $s'=v^2s$. These charts give the line bundle $\mathcal O_{\mathbb P^1}(-2)$ near its zero section. The quotient is therefore smooth and resolves the double point, with the stated exceptional curve. Doing this at all fixed points gives the compact surface $X$.

Choose a flat K\"ahler form on $A$ that locally equals $\ddc\rho$, where $\rho=|z_1|^2+|z_2|^2$, on disjoint balls centered at the fixed points. Away from those points it descends to $Y$. On the resolution chart,
\[
 \rho=|s|(1+|v|^2),\qquad
 \rho^2=|s|^2(1+|v|^2)^2.
\]
Choose a smooth function $F$ on $[0,+\infty)$ with $F(\rho)=c\rho^2$ near zero, $c>0$, $F'>0$ on $(0,+\infty)$, $F''\geq0$, and $F(\rho)=\rho-C$ for $\rho\geq\delta$, for a sufficiently small $\delta>0$. Such a function is obtained by choosing $F'$ nondecreasing, linear near zero and equal to one past $\delta$. Replace the flat form in each ball by $\ddc F(\rho)$. This agrees with the flat form near the boundary and extends smoothly across the exceptional line. The resulting closed form $\eta$ is positive off the exceptional set; along it, $\eta$ is positive in the normal direction and zero in the tangent direction.

To restore tangent positivity, choose a smooth cutoff $\chi$ equal to one for $\rho\leq\delta$ and zero for $\rho\geq2\delta$. For $s\ne0$, set
\[
 q_v=\chi(\rho)\log\rho-\log|s|.
\]
Near $s=0$ this equals $\log(1+|v|^2)$ and extends smoothly. On the overlap, $q_{v'}-q_v=-2\log|v|$, so their $\ddc$ forms agree. They define a smooth closed form $R$ supported in the resolved ball, positive in the tangent direction along the exceptional line and semipositive where $\chi=1$. The form $\eta+\varepsilon R$ is positive there, including along the line. On the compact transition annulus $\eta$ is already positive, so it remains positive for sufficiently small $\varepsilon>0$. Outside the ball $R=0$. Summing these corrections over the disjoint balls produces a K\"ahler form on $X$.
\end{proof}

Fix now $a\in\mathbb R\setminus\mathbb Q$ and use the lattice generated by
\[
 \lambda_1=(1,0),\quad \lambda_2=(a,1),\quad
 \lambda_3=(\mathrm i,0),\quad \lambda_4=(0,\mathrm i).
\]
Write $z_j=x_j+\mathrm i y_j$. The coordinate $y_1\bmod\mathbb Z$ defines a map from $A$ to the circle. Choose a smooth even function $h:\mathbb R/\mathbb Z\to[0,+\infty)$ with integral one and vanishing near $\frac12\mathbb Z$. Such functions are obtained from even pairs of nonnegative bumps away from that set. Put
\[
 \theta_h=h(y_1)\,\mathrm dx_1\wedge\mathrm dy_1.
\]

\Needspace{14\baselineskip}
\begin{proposition}\label{prop:kummer-nonrigid-irrational}\index{Nef class!irrational nonrigid class}\index{Semipositive form!quotient resolution}
The forms $\theta_h$ descend and extend by zero to smooth semipositive closed forms $T_h$ on $X$. They all represent one nonzero nef Bott--Chern class $\alpha$ with $\alpha^2=0$. This class is nonrigid, and its image in real de Rham cohomology has an irrational ray. Each $T_h$ vanishes on a neighborhood of the exceptional set.
\end{proposition}
\begin{proof}
The form $\theta_h$ is closed, because $\mathrm dh(y_1)$ is proportional to $\mathrm dy_1$, and is semipositive. Evenness makes it invariant under the involution. At every fixed point, $y_1$ belongs to $\frac12\mathbb Z$ modulo $\mathbb Z$, so the form is zero on a neighborhood of that point. Its descent away from the singular points therefore extends by zero on their resolved neighborhoods. This proves smoothness and the asserted vanishing.

For two admissible functions $h_1,h_2$, put $k=h_1-h_2$. It is even, has mean zero and vanishes near $\frac12\mathbb Z$. Choose $f$ with
\[
 f'(y)=4\pi\int_0^y k(s)\,\mathrm ds,
 \qquad f''=4\pi k.
\]
The derivative is periodic and odd. Its mean is zero, so $f$ can be chosen periodic and even. Evenness and the zero mean of $k$ give $\int_0^{1/2}k=0$; hence $f'=0$ near $\frac12\mathbb Z$. The function $f(y_1)$ descends and extends constantly near each exceptional curve to a smooth function $G$ on $X$. With the book's normalization,
\[
 \ddc G=T_{h_1}-T_{h_2},
 \qquad
 \ddc f(y_1)=\frac{f''(y_1)}{4\pi}
                    \,\mathrm dx_1\wedge\mathrm dy_1
 \quad\text{on }A.
\]
Thus the forms have the same Bott--Chern class. Distinct admissible functions give distinct forms, proving nonrigidity. They are semipositive and have square zero. By \cref{lem:kummer-resolution-kahler}, adding any positive multiple of a K\"ahler form gives a K\"ahler representative of a nearby class; therefore $\alpha$ is nef.

The periods of $\theta_h$ on the real two-tori generated by $(\lambda_1,\lambda_3)$ and $(\lambda_2,\lambda_3)$ are respectively $1$ and $a$, by the mean-one condition. Translate both cycles by $(0,\mathrm i/4)$. They have $y_2=1/4$ modulo $\mathbb Z$ and avoid the sixteen fixed points; their periods are unchanged because the form is closed. They lie in the unramified torus cover of the complement of the exceptional curves. The first period proves that $\alpha\ne0$. If a nonzero multiple $b\alpha$ had rational de Rham class, pulling back and evaluating on these cycles would make both $b$ and $ba$ rational, contradicting $a\notin\mathbb Q$. Finally, the zero square prevents a K\"ahler representative of $\alpha$ itself by Stokes' theorem, so it is a boundary class.
\end{proof}
